\documentclass[12pt]{article}

% ===================== PAQUETES =====================
\usepackage[utf8]{inputenc}
\usepackage[spanish]{babel}
\usepackage[margin=2.5cm]{geometry}
\usepackage{graphicx}
\usepackage{xcolor}
\usepackage{mathptmx}
\usepackage{microtype}
\usepackage{colortbl}
\usepackage{enumitem}
\usepackage{titlesec}
\usepackage{fancyhdr}
\usepackage{float}
\usepackage[hidelinks]{hyperref}
\usepackage{setspace}
\usepackage{longtable}

% ===================== COLORES =====================
\definecolor{ipn-guinda}{RGB}{135, 30, 48}
\definecolor{escom-azul}{RGB}{0, 102, 160}

% ===================== ESPACIADO DE PÁRRAFO =====================
\setlength{\parindent}{0pt}
\setlength{\parskip}{0.6em}

% ===================== ESTILO DE SECCIONES =====================
\titleformat{\section}
  {\color{ipn-guinda}\Large\bfseries}{\thesection}{0.5em}{}
\titleformat{\subsection}
  {\color{escom-azul}\large\bfseries}{\thesubsection}{0.5em}{}
\titlespacing*{\section}{0pt}{1.4em}{0.8em}
\titlespacing*{\subsection}{0pt}{1em}{0.5em}

% ===================== ENCABEZADO / PIE DE PÁGINA =====================
\pagestyle{fancy}
\setlength{\headheight}{14pt}
\fancyhf{}
\renewcommand{\headrulewidth}{0.4pt}
\fancyhead[L]{\small Analítica y Visualización de Datos}
\fancyhead[R]{\small Proyecto Final}
\fancyfoot[C]{\small\thepage}

\begin{document}

% ===================== PORTADA =====================
\begin{titlepage}
    \thispagestyle{empty}
    \centering
    \begin{minipage}{0.3\textwidth}
        \centering
        \includegraphics[width=0.55\textwidth]{ipnescudo.png}
    \end{minipage}
    \hfill
    \begin{minipage}{0.3\textwidth}
        \centering
        \includegraphics[width=0.55\textwidth]{escudoESCOM.png}
    \end{minipage}

    \vspace{1.5cm}
    {\color{ipn-guinda} \LARGE \textbf{INSTITUTO POLITÉCNICO NACIONAL}} \\
    \vspace{0.2cm}
    {\color{escom-azul} \large \textbf{ESCUELA SUPERIOR DE CÓMPUTO}}

    \vspace{1.5cm}
    \rule{\textwidth}{0.8pt} \\
    \vspace{0.8cm}
    {\LARGE \textbf{Proyecto Final}} \\
    \vspace{0.3cm}
    {\large Analítica y Visualización de Datos}
    \vspace{0.5cm}
    \rule{\textwidth}{0.8pt}

    \vspace{1.5cm}
    \textbf{\Large Profesor:} \\
    \vspace{0.2cm}
    Alejandro López Gómez

    \vspace{2cm}
    \textbf{\Large Alumnos:} \\
    \vspace{0.2cm}
    García Juárez Dayan \\
    Munguía Aguilera Cesar Raúl

    \vspace{1.2cm}
    \textbf{\Large Grupo:} \\
    \vspace{0.2cm}
    5AM1

    \vfill
    {\large Ciudad de México, \today}
\end{titlepage}

% ===================== ÍNDICE =====================
\onehalfspacing
\tableofcontents
\newpage

% ===================== INTRODUCCIÓN =====================
\section*{Introducción}
\addcontentsline{toc}{section}{Introducción}

En un entorno empresarial cada vez más competitivo, la capacidad de transformar datos en información útil para la toma de decisiones constituye una ventaja relevante. En este contexto, el presente proyecto analiza el comportamiento de consumo de la base de clientes de la organización con el fin de identificar patrones de compra, describir perfiles de alto valor y aportar evidencia para orientar decisiones comerciales.

El estudio parte de un proceso de preparación, limpieza y análisis de la información, seguido por la construcción de indicadores de valor, actividad y antigüedad. El alcance es principalmente descriptivo y exploratorio: las asociaciones encontradas permiten formular hipótesis comerciales, pero no demuestran relaciones causales ni permiten atribuir por sí solas las decisiones de compra a una variable específica.

Se presta especial atención a los clientes de alto valor y a aquellos que presentan gasto elevado desde etapas tempranas de su relación con la empresa. La identificación de estos grupos puede servir como base para diseñar, en una etapa posterior, estrategias de comunicación, retención y fidelización.

Los resultados sugieren que las variables económicas y transaccionales aportan señales útiles para distinguir perfiles de consumo. El dashboard, las conclusiones finales y las recomendaciones comerciales se dejan pendientes hasta definir su alcance y validar las visualizaciones definitivas.

\newpage
\section*{Contexto del Dataset y Problema de Negocio}
\addcontentsline{toc}{section}{Contexto del Dataset y Problema de Negocio}

El archivo original utilizado en el proyecto se denomina \textit{ml\_project1\_data}. El conjunto de datos se asocia con el \textbf{iFood CRM Data Analyst Case}, un caso práctico empleado para evaluar candidatos de análisis de datos en el equipo \textit{iFood Brain}. Una fuente pública trazable es el repositorio \texttt{ifood-data-business-analyst-test}, cuyo README señala que el caso se utilizaba en la contratación de Data Analysts para esa área.\footnote{Repositorio \texttt{nailson/ifood-data-business-analyst-test}, README: \url{https://github.com/nailson/ifood-data-business-analyst-test}.}

La institución relacionada es \textbf{iFood}, una empresa brasileña de tecnología enfocada en servicios de entrega en línea. Su sitio institucional la presenta como una plataforma que conecta clientes, establecimientos y repartidores.\footnote{Sitio institucional de iFood: \url{https://institucional.ifood.com.br/}. Sitio comercial de iFood: \url{https://www.ifood.com.br/}.}

El caso describe una empresa del sector retail/alimentos. Las categorías de producto registradas son vinos, frutas, carnes, pescados, dulces y productos \textit{gold}. Los canales de compra analizados incluyen tienda física, catálogo y sitio web.

El problema central consiste en estudiar el desempeño de una campaña de marketing directo. Se registraron \textbf{2{,}240 clientes}; \texttt{Response} identifica a quienes aceptaron la campaña más reciente. En total se observaron \textbf{334 respuestas positivas}, equivalentes a una tasa de aceptación de \textbf{14.91\%}.

\begin{table}[h!]
\centering
\small
\begin{tabular}{|l|r|}
\hline
\rowcolor{escom-azul!20}
\textbf{Concepto} & \textbf{Valor} \\
\hline
Clientes registrados & 2{,}240 \\
\hline
Respuestas positivas & 334 \\
\hline
Costo total de contacto & 6{,}720 MU \\
\hline
Ingreso asociado a respuestas & 3{,}674 MU \\
\hline
Balance simplificado & $-3{,}046$ MU \\
\hline
Tasa de aceptación & 14.91\% \\
\hline
\end{tabular}
\caption{Resumen descriptivo y financiero de la campaña del caso iFood CRM}
\label{tab:piloto_ifood}
\end{table}

No existe una fuente oficial que especifique la divisa de los montos. Por ello, se usa \textbf{MU} como abreviatura de unidades monetarias genéricas. Las variables \texttt{Z\_CostContact = 3} y \texttt{Z\_Revenue = 11} permiten reproducir los cálculos del caso:

\[
2240 \times 3 = 6720 \text{ MU}
\]
\[
334 \times 11 = 3674 \text{ MU}
\]

El balance mostrado corresponde a la diferencia simple entre ingreso y costo, $3674-6720=-3046$ MU. No debe interpretarse como utilidad contable completa, ya que el dataset no contiene todos los costos e ingresos de la operación.


% ===================== OBJETIVOS =====================
\section*{Objetivos del Proyecto}
\addcontentsline{toc}{section}{Objetivos del Proyecto}

\subsection*{Objetivo General}
\addcontentsline{toc}{subsection}{Objetivo General}

Analizar el comportamiento y los patrones de consumo de la base de clientes mediante técnicas de análisis de datos, con el propósito de identificar segmentos estratégicos y generar evidencia que pueda apoyar la asignación de recursos en futuras campañas de marketing.

\subsection*{Objetivos Específicos}
\addcontentsline{toc}{subsection}{Objetivos Específicos}

\begin{itemize}[leftmargin=1.5em, itemsep=0.4em]
    \item \textbf{Limpieza y consolidación de datos:} diagnosticar y tratar inconsistencias, registros no válidos y valores faltantes para disponer de una base coherente para el análisis.
    \item \textbf{Segmentación de clientes:} clasificar la cartera mediante indicadores de valor, estructura familiar y antigüedad para comparar perfiles de consumo.
    \item \textbf{Análisis de campañas:} describir el desempeño histórico de la campaña más reciente mediante su tasa de aceptación y la comparación de perfiles de respuesta.
    \item \textbf{Preparación para el dashboard:} generar desde el notebook un archivo analítico derivado que sirva como fuente única para el futuro tablero. El diseño e implementación del dashboard se mantienen pendientes.
\end{itemize}

\newpage

% ===================== METODOLOGÍA =====================
\section*{Metodología}
\addcontentsline{toc}{section}{Metodología}

El proyecto se desarrolló mediante un proceso iterativo organizado en cuatro etapas. Las tres primeras se encuentran implementadas en el notebook; la cuarta queda pendiente hasta definir el dashboard.

\subsection*{Fase 1: Análisis Exploratorio de Datos (EDA) y Limpieza}
\addcontentsline{toc}{subsection}{Fase 1: EDA y Limpieza}

Se inspeccionaron dimensiones, tipos de datos, valores nulos, duplicados, constantes y rangos. Después se convirtió \texttt{Dt\_Customer} a fecha, se imputaron los valores faltantes de \texttt{Income} con su mediana y se eliminaron únicamente tres registros con edades no plausibles. Las variables monetarias y de gasto se conservaron, incluso cuando presentaban valores extremos, para no eliminar clientes potencialmente válidos de alto valor.

\subsection*{Fase 2: Ingeniería de Variables}
\addcontentsline{toc}{subsection}{Fase 2: Ingeniería de Variables}

Se crearon \texttt{Age}, \texttt{Customer\_Tenure\_Days}, \texttt{Total\_Spending}, \texttt{Total\_Purchases}, \texttt{Average\_Ticket}, \texttt{Total\_Dependents}, \texttt{Accepted\_Campaigns\_Total}, \texttt{Value\_Segment}, \texttt{Family\_Segment} y \texttt{Tenure\_Group}. \texttt{Total\_Purchases} suma las compras por web, catálogo y tienda; no incluye \texttt{NumDealsPurchases}, porque esta última variable describe compras con descuento y podría solaparse con los canales.

\subsection*{Fase 3: Análisis de Relaciones y Segmentación}
\addcontentsline{toc}{subsection}{Fase 3: Análisis de Relaciones y Segmentación}

\begin{itemize}[leftmargin=1.5em, itemsep=0.4em]
    \item \textbf{Estadística descriptiva y análisis univariado:} caracterización de distribuciones, rangos y valores extremos.
    \item \textbf{Correlación:} medición de asociaciones lineales entre variables numéricas. Estas asociaciones no se interpretan como causalidad.
    \item \textbf{Segmentación:} comparación de clientes por valor, antigüedad y estructura familiar reportada.
    \item \textbf{Campaña:} comparación descriptiva entre clientes que aceptaron y no aceptaron la campaña más reciente.
    \item \textbf{PCA:} representación exploratoria de seis variables estandarizadas en dos componentes principales.
\end{itemize}

\subsection*{Fase 4: Visualización y Dashboarding}
\addcontentsline{toc}{subsection}{Fase 4: Visualización y Dashboarding}

\textit{Pendiente.} El notebook genera \texttt{ifood\_dashboard.csv} como fuente derivada para el tablero, pero el diseño, las métricas visibles y la tecnología del dashboard se definirán posteriormente.

% ===================== DATOS =====================
\section*{Descripción del Dataset}
\addcontentsline{toc}{section}{Descripción de los Datos}

Para este análisis se utilizó el conjunto de datos original \textit{ml\_project1\_data}, compuesto por \textbf{2{,}240 registros y 29 variables}. En la estructura local del proyecto se conserva como \texttt{data/ifood.csv}; este cambio de nombre no altera su procedencia ni su contenido. Después de la limpieza y la ingeniería de variables, el notebook genera una tabla analítica de \textbf{2{,}237 registros y 39 columnas} para alimentar el futuro dashboard.

\subsection*{Estructura del Dataset}
\addcontentsline{toc}{subsection}{Estructura del Dataset}

\begin{itemize}[leftmargin=1.5em, itemsep=0.4em]
    \item \textbf{Variables sociodemográficas:} año de nacimiento, educación, estado civil, ingreso y dependientes en el hogar.
    \item \textbf{Variables transaccionales:} gasto por categoría y número de compras por canal o con descuento.
    \item \textbf{Variables de interacción y campañas:} fecha de alta, recencia, visitas web, quejas y respuestas a campañas.
    \item \textbf{Variables derivadas:} edad histórica, antigüedad, gasto total, compras totales, ticket promedio y segmentos analíticos.
\end{itemize}

\subsection*{Diccionario de Datos}
\addcontentsline{toc}{subsection}{Diccionario de Datos}

\scriptsize
\begin{longtable}{|p{3.1cm}|p{4.6cm}|p{6.2cm}|}
\caption{Diccionario de las 29 variables originales de \textit{ml\_project1\_data}}
\label{tab:diccionario}\\
\hline
\rowcolor{escom-azul!20}
\textbf{Variable} & \textbf{Significado} & \textbf{Contexto útil} \\
\hline
\endfirsthead
\hline
\rowcolor{escom-azul!20}
\textbf{Variable} & \textbf{Significado} & \textbf{Contexto útil} \\
\hline
\endhead
\texttt{ID} & Identificador único del cliente & Distingue registros; se excluye como predictor numérico. \\
\hline
\texttt{Year\_Birth} & Año de nacimiento & Permite calcular la edad usando una fecha de referencia coherente con el periodo del dataset. \\
\hline
\texttt{Education} & Nivel educativo & Variable sociodemográfica categórica. \\
\hline
\texttt{Marital\_Status} & Estado civil reportado & Aporta contexto sobre la estructura declarada del hogar. \\
\hline
\texttt{Income} & Ingreso anual del hogar & Monto en unidades monetarias no especificadas. \\
\hline
\texttt{Kidhome} & Número de niños pequeños en casa & Parte del total de dependientes usado en la segmentación familiar. \\
\hline
\texttt{Teenhome} & Número de adolescentes en casa & Complementa el total de dependientes del hogar. \\
\hline
\texttt{Dt\_Customer} & Fecha de alta como cliente & Permite calcular la antigüedad con respecto al 29 de junio de 2014. \\
\hline
\texttt{Recency} & Días desde la última compra & Un valor menor indica una compra más reciente. \\
\hline
\texttt{MntWines} & Monto gastado en vinos en dos años & Componente del gasto total. \\
\hline
\texttt{MntFruits} & Monto gastado en frutas en dos años & Componente del gasto total. \\
\hline
\texttt{MntMeatProducts} & Monto gastado en carnes en dos años & Componente del gasto total. \\
\hline
\texttt{MntFishProducts} & Monto gastado en pescado en dos años & Componente del gasto total. \\
\hline
\texttt{MntSweetProducts} & Monto gastado en dulces en dos años & Componente del gasto total. \\
\hline
\texttt{MntGoldProds} & Monto gastado en productos \textit{gold} en dos años & Componente del gasto total. \\
\hline
\texttt{NumDealsPurchases} & Compras realizadas con descuento & Mide uso de promociones; no se suma a los canales para evitar posible doble conteo. \\
\hline
\texttt{NumWebPurchases} & Compras por sitio web & Componente de \texttt{Total\_Purchases}. \\
\hline
\texttt{NumCatalogPurchases} & Compras por catálogo & Componente de \texttt{Total\_Purchases}. \\
\hline
\texttt{NumStorePurchases} & Compras en tienda física & Componente de \texttt{Total\_Purchases}. \\
\hline
\texttt{NumWebVisitsMonth} & Visitas al sitio web en el último mes & Mide actividad de navegación; no equivale por sí misma a intención ni a conversión. \\
\hline
\texttt{AcceptedCmp1} & Aceptó la campaña 1 & Indicador binario de respuesta histórica. \\
\hline
\texttt{AcceptedCmp2} & Aceptó la campaña 2 & Indicador binario de respuesta histórica. \\
\hline
\texttt{AcceptedCmp3} & Aceptó la campaña 3 & Indicador binario de respuesta histórica. \\
\hline
\texttt{AcceptedCmp4} & Aceptó la campaña 4 & Indicador binario de respuesta histórica. \\
\hline
\texttt{AcceptedCmp5} & Aceptó la campaña 5 & Indicador binario de respuesta histórica. \\
\hline
\texttt{Response} & Aceptó la campaña más reciente & Indicador binario usado para comparar perfiles de respuesta. \\
\hline
\texttt{Complain} & Registró una queja en los últimos dos años & Aporta contexto sobre la experiencia del cliente. \\
\hline
\texttt{Z\_CostContact} & Costo unitario de contacto & Es constante e igual a 3; se conserva para los cálculos financieros. \\
\hline
\texttt{Z\_Revenue} & Ingreso unitario asociado a una respuesta & Es constante e igual a 11; se conserva para los cálculos financieros. \\
\hline
\end{longtable}
\normalsize

\newpage

% ===================== EDA =====================
\section*{Análisis Exploratorio de Datos (EDA)}
\addcontentsline{toc}{section}{Análisis Exploratorio de Datos (EDA)}

El análisis exploratorio se estructuró en seis etapas: diagnóstico de integridad, estadística descriptiva, análisis univariado, análisis bivariado, análisis por segmentos y reducción de dimensionalidad. Los resultados siguientes corresponden a la versión ejecutada del notebook con 2{,}237 registros limpios.

\subsection*{1. Diagnóstico de Integridad y Calidad de Datos}
\addcontentsline{toc}{subsection}{1. Diagnóstico de Integridad y Calidad de Datos}

El dataset original \textit{ml\_project1\_data} contiene 2{,}240 registros y 29 variables. La única columna original con faltantes es \texttt{Income}, con 24 casos (1.07\%). El dataset no permite determinar por qué faltan esos valores, por lo que no se atribuye su ausencia a una causa concreta. No se detectaron registros duplicados.

\begin{figure}[h!]
    \centering
    \includegraphics[width=0.85\textwidth]{datos_faltantes.png}
    \caption{Mapa de calor de valores faltantes en el dataset original}
    \label{fig:datos_faltantes}
\end{figure}

Los 24 valores faltantes se imputaron con la mediana de \texttt{Income}, igual a 51{,}373 MU. Esta decisión conserva el tamaño de la muestra y es menos sensible que la media al valor extremo de 666{,}666 MU.

\subsection*{2. Estadística Descriptiva}
\addcontentsline{toc}{subsection}{2. Estadística Descriptiva}

\begin{itemize}[leftmargin=1.5em, itemsep=0.4em]
    \item \textbf{Gasto total:} la media de \texttt{Total\_Spending} (605.74 MU) supera la mediana (396 MU), lo que refleja una distribución asimétrica hacia valores altos.
    \item \textbf{Ingresos:} \texttt{Income} incluye un máximo de 666{,}666 MU. El valor se conserva en la base limpia; únicamente se limita de manera temporal al preparar el PCA.
    \item \textbf{Recencia:} \texttt{Recency} abarca de 0 a 99 días, con media y mediana próximas a 49 y una asimetría cercana a cero ($-0.003$), por lo que su distribución presenta poca asimetría.
\end{itemize}

\begin{table}[h!]
\centering
\small
\begin{tabular}{|l|r|r|r|r|r|}
\hline
\rowcolor{escom-azul!20}
\textbf{Variable} & \textbf{Media} & \textbf{Mediana} & \textbf{Desv. Est.} & \textbf{Mín.} & \textbf{Máx.} \\
\hline
\texttt{Income}            & 52{,}227.32 & 51{,}373.00 & 25{,}043.27 & 1{,}730.00 & 666{,}666.00 \\
\hline
\texttt{Total\_Spending}   & 605.74      & 396.00      & 601.84      & 5.00       & 2{,}525.00 \\
\hline
\texttt{Age}               & 45.10       & 44.00       & 11.70       & 18.00      & 74.00 \\
\hline
\texttt{Recency}           & 49.10       & 49.00       & 28.96       & 0.00       & 99.00 \\
\hline
\texttt{NumWebVisitsMonth} & 5.32        & 6.00        & 2.43        & 0.00       & 20.00 \\
\hline
\end{tabular}
\caption{Medidas descriptivas calculadas sobre los 2{,}237 registros limpios}
\label{tab:estadisticas}
\end{table}

\subsection*{3. Análisis Univariado}
\addcontentsline{toc}{subsection}{3. Análisis Univariado}

El análisis univariado caracterizó la distribución individual de las variables y permitió identificar rangos, asimetrías y valores extremos. La mayoría de los ingresos se concentra muy por debajo del máximo de 666{,}666 MU. Este valor extremo se conservó porque no existe evidencia suficiente para declararlo un error de captura.

\begin{figure}[H]
    \centering
    \includegraphics[width=\textwidth]{analisis_univariado.png}
    \caption{Distribuciones de variables numéricas seleccionadas}
    \label{fig:univariado}
\end{figure}

\texttt{NumWebVisitsMonth} tiene media 5.32 y mediana 6 visitas. Esta variable describe frecuencia de navegación, pero por sí sola no permite medir la calidad del involucramiento digital ni distinguir visitas con intención de compra.

\textbf{Composición agregada del gasto por categoría:} al dividir la suma de cada categoría entre el gasto total de toda la muestra limpia, vinos representa 50.19\%, carnes 27.56\%, productos \textit{gold} 7.26\%, pescado 6.19\%, dulces 4.47\% y frutas 4.34\%. Estas cifras describen la participación agregada de cada categoría y no el promedio de porcentajes individuales por cliente.

\begin{figure}[H]
    \centering
    \includegraphics[width=0.4\textwidth]{mntwines.png}
    \caption{Participación agregada del gasto por categoría de producto}
    \label{fig:composicion_gasto}
\end{figure}

\textbf{Distribución de edades y gasto total:} usando como referencia el 29 de junio de 2014, la edad varía entre 18 y 74 años. La mediana es 44 años y el 50\% central se encuentra entre 37 y 55 años. El gasto total presenta una cola hacia valores altos.

\begin{figure}[H]
    \centering
    \includegraphics[width=0.8\textwidth]{distribucionedad.png}
    \caption{Distribución de edad histórica y gasto total de los clientes}
    \label{fig:distribucion_edad}
\end{figure}

El percentil 90 de \texttt{Total\_Spending} es 1{,}536 MU. Los 225 clientes con gasto igual o superior a ese umbral concentran aproximadamente 30.29\% del gasto y presentan una tasa descriptiva de respuesta de 42.22\%. Esta asociación identifica un segmento relevante, pero no demuestra que el gasto cause la aceptación de la campaña.

\textbf{Detección de valores extremos mediante boxplots:} las variables monetarias y de gasto presentan puntos por encima de los bigotes. Los boxplots se utilizan como herramienta descriptiva; los puntos no se eliminan automáticamente, ya que pueden corresponder a clientes reales de alto valor. La única exclusión de registros aplicada en la limpieza se debió al criterio de edad válida.

\begin{figure}[H]
    \centering
    \includegraphics[width=0.7\textwidth]{outiliers.png}
    \caption{Detección descriptiva de valores extremos mediante boxplots}
    \label{fig:outliers}
\end{figure}

\subsection*{4. Análisis Bivariado y Correlación}
\addcontentsline{toc}{subsection}{4. Análisis Bivariado y Correlación}

El análisis bivariado exploró asociaciones lineales entre pares de variables. Todas las correlaciones reportadas son descriptivas y no implican causalidad.

\textbf{Ingreso y gasto:} \texttt{Income} se correlaciona con \texttt{Total\_Spending} en $r=0.665$, con \texttt{MntWines} en $r=0.577$, con \texttt{MntMeatProducts} en $r=0.578$ y con \texttt{MntGoldProds} en $r=0.321$. Las asociaciones son positivas, pero su intensidad varía por categoría.

\begin{figure}[H]
    \centering
    \includegraphics[width=0.85\textwidth]{correlacion.png}
    \caption{Matriz de correlación entre variables numéricas seleccionadas}
    \label{fig:correlacion}
\end{figure}

Los hallazgos principales son:

\begin{itemize}[leftmargin=1.5em, itemsep=0.4em]
    \item \textbf{Ingreso y gasto:} las asociaciones más altas de \texttt{Income} entre las variables seleccionadas se observan con el gasto total, vinos y carnes.
    \item \textbf{Canales de compra:} la correlación entre \texttt{NumCatalogPurchases} y \texttt{NumWebPurchases} es $r=0.378$; es positiva, pero no debe describirse como fuerte.
    \item \textbf{Visitas web:} \texttt{NumWebVisitsMonth} presenta correlaciones de $r=-0.549$ con \texttt{Income} y $r=-0.500$ con \texttt{Total\_Spending}. El dataset no permite concluir que estas personas naveguen sin intención de comprar.
    \item \textbf{Dependientes:} \texttt{Kidhome} presenta una correlación de $r=-0.557$ con \texttt{Total\_Spending}, mientras que \texttt{Teenhome} presenta $r=-0.138$. El patrón es distinto para cada tipo de dependiente.
    \item \textbf{Variables constantes:} \texttt{Z\_CostContact} y \texttt{Z\_Revenue} se conservan para los cálculos financieros, pero se excluyen de la matriz porque no tienen varianza.
\end{itemize}

\textbf{Respuesta a la campaña más reciente y canales de compra:} las medianas de compras para clientes sin respuesta fueron 3 por web, 1 por catálogo y 5 en tienda. Para clientes con respuesta positiva fueron 5 por web, 4 por catálogo y 6 en tienda. La comparación corresponde a \texttt{Response}, no a la suma de campañas previas.

\begin{figure}[H]
    \centering
    \includegraphics[width=0.7\textwidth]{campañas.png}
    \caption{Compras por canal según la respuesta a la campaña más reciente (\texttt{Response})}
    \label{fig:campanas}
\end{figure}

\subsection*{5. Análisis por Segmentos}
\addcontentsline{toc}{subsection}{5. Análisis por Segmentos}

\textbf{Antigüedad del cliente:} \texttt{Customer\_Tenure\_Days} se calculó desde \texttt{Dt\_Customer} hasta el 29 de junio de 2014. Después, la base se dividió por cuantiles en cuatro grupos: Nuevo, Reciente, Estable y Leal.

\begin{figure}[H]
    \centering
    \includegraphics[width=0.85\textwidth]{gastos.png}
    \caption{Gasto total por grupo de antigüedad del cliente}
    \label{fig:gastos_antiguedad}
\end{figure}

\begin{itemize}[leftmargin=1.5em, itemsep=0.4em]
    \item \textbf{Medianas:} el gasto mediano es 169.5 MU en Nuevo, 325 MU en Reciente, 457.5 MU en Estable y 583.5 MU en Leal.
    \item \textbf{Asociación débil:} la correlación de Pearson entre antigüedad y gasto total es $r=0.159$. Por tanto, el patrón creciente de medianas no implica que la permanencia sea el único factor asociado al gasto ni que exista una relación causal fuerte.
    \item \textbf{Valores máximos:} los máximos observados son 2{,}525 MU en Nuevo, 2{,}524 MU en Reciente, 2{,}440 MU en Estable y 2{,}352 MU en Leal. Existen clientes de gasto alto en todos los segmentos.
    \item \textbf{Alto valor temprano:} se identificaron 38 clientes de los grupos Nuevo o Reciente por encima del percentil 95 del gasto total, cuyo umbral es 1{,}767.2 MU.
\end{itemize}

No se afirma que el grupo Leal sea simultáneamente el más disperso y el más predecible; ambas ideas requieren métricas específicas de dispersión y estabilidad temporal que no forman parte de este análisis transversal.

\subsection*{Estado civil y estructura familiar}
\addcontentsline{toc}{subsection}{Estado civil y estructura familiar}

La segmentación combina \texttt{Total\_Dependents} con el estado civil reportado y distingue tres grupos. Esta clasificación describe los datos disponibles; no verifica la estructura real del hogar más allá de las variables registradas.

\begin{table}[H]
\centering
\small
\begin{tabular}{|l|r|r|r|}
\hline
\rowcolor{escom-azul!20}
\textbf{Segmento} & \textbf{Clientes} & \textbf{Mediana (MU)} & \textbf{Media (MU)} \\
\hline
Con dependientes en pareja & 1{,}060 & 215 & 413.63 \\
\hline
Con dependientes sin pareja reportada & 540 & 180 & 394.09 \\
\hline
Sin dependientes & 637 & 1{,}189 & 1{,}104.86 \\
\hline
\end{tabular}
\caption{Gasto total por segmento familiar reportado}
\label{tab:segmento_familiar}
\end{table}

\begin{figure}[H]
    \centering
    \includegraphics[width=0.80\textwidth]{edocivil.png}
    \caption{Gasto total por segmento familiar reportado}
    \label{fig:edocivil}
\end{figure}

Entre los dos grupos con dependientes, las medianas son relativamente cercanas (215 y 180 MU). Sin embargo, el grupo sin dependientes presenta una mediana de 1{,}189 MU, por lo que no es correcto generalizar que la estructura familiar carece de relación con el gasto. Las diferencias son descriptivas y pueden estar asociadas también con ingreso, edad u otras variables; no se realizó un modelo multivariado para aislar esos efectos.

\subsection*{6. Reducción de Dimensionalidad (PCA)}
\addcontentsline{toc}{subsection}{6. Reducción de Dimensionalidad (PCA)}

El PCA se aplicó a seis variables: \texttt{Income}, \texttt{Total\_Spending}, \texttt{Age}, \texttt{Customer\_Tenure\_Days}, \texttt{Recency} y \texttt{NumWebVisitsMonth}. Para reducir la influencia de extremos exclusivamente en este procedimiento, cada variable se limitó temporalmente a sus percentiles 1 y 99 y después se estandarizó. Este recorte no modifica los valores conservados en la tabla limpia ni en el archivo destinado al dashboard.

El primer componente explica \textbf{39.72\%} de la varianza y el segundo \textbf{19.34\%}; juntos explican \textbf{59.06\%}. En consecuencia, la proyección en dos dimensiones resume una parte mayoritaria, pero conserva fuera del plano aproximadamente 40.94\% de la variabilidad.

\begin{figure}[H]
    \centering
    \includegraphics[width=0.80\textwidth]{pca.png}
    \caption{Proyección de los clientes sobre los dos primeros componentes principales}
    \label{fig:pca}
\end{figure}

La proyección permite explorar gradientes y observaciones alejadas, pero no demuestra por sí sola la existencia de grupos discretos. Además, el signo de cada componente es arbitrario; por ello, no se asigna automáticamente el lado derecho o izquierdo del gráfico a clientes de alto valor sin revisar las cargas del componente.

% ===================== SECCIONES PENDIENTES =====================
\newpage
\section*{Dashboard Interactivo}
\addcontentsline{toc}{section}{Dashboard Interactivo}

El análisis se complementó con un dashboard interactivo desarrollado en \textbf{Streamlit}. Su propósito es facilitar la consulta de los resultados del EDA y apoyar la comparación histórica de clientes, canales, productos y campañas sin incorporar un modelo predictivo. La aplicación permite modificar filtros y recalcular indicadores sobre subconjuntos de la muestra, pero no estima la respuesta futura de un cliente ni atribuye causalidad a las asociaciones observadas.

\subsection*{Arquitectura y fuente de datos}
\addcontentsline{toc}{subsection}{Arquitectura y fuente de datos}

La aplicación se organiza mediante un archivo principal \texttt{app.py} y siete módulos dentro de la carpeta \texttt{pages/}. El notebook \texttt{Analitica.ipynb} se mantiene como la única fuente de verdad del proyecto: realiza la limpieza, crea las variables derivadas y exporta el archivo \texttt{data/ifood\_dashboard.csv}. El dashboard se limita a leer dicho archivo, aplicar filtros, generar agregaciones y presentar las visualizaciones, evitando duplicar reglas de preparación de datos dentro de la interfaz.

La tabla analítica exportada contiene \textbf{2{,}237 clientes válidos y 61 columnas}. Además, conserva como metadatos los \textbf{2{,}240 contactos} y las \textbf{334 respuestas positivas} de la campaña original. Esta separación permite utilizar la muestra limpia para el análisis de perfiles sin alterar el resumen financiero histórico del piloto.

\subsection*{Estructura de navegación}
\addcontentsline{toc}{subsection}{Estructura de navegación}

El dashboard sigue un flujo narrativo de siete páginas:

\begin{enumerate}[leftmargin=1.8em, itemsep=0.35em]
    \item \textbf{Resumen ejecutivo:} presenta clientes analizados, gasto promedio y mediano, tasa de aceptación y resultado financiero de la campaña original.
    \item \textbf{Perfil de clientes:} describe edad, educación, estado civil, ingreso y composición de dependientes.
    \item \textbf{Valor y consumo:} analiza la distribución de \texttt{Total\_Spending} y la participación de vinos, carnes, frutas, pescado, dulces y productos \textit{gold}.
    \item \textbf{Canales de compra:} compara web, catálogo, tienda, visitas digitales y compras con descuento.
    \item \textbf{Campañas:} muestra la aceptación de las campañas 1 a 5 y de la campaña más reciente, junto con comparaciones por ingreso, gasto y canal.
    \item \textbf{Segmentos estratégicos:} resume cohortes de interés mediante su tamaño, gasto, aceptación, balance, canal y categoría principal.
    \item \textbf{Centro de decisiones:} funciona como un simulador retrospectivo que compara el contacto masivo con una selección definida mediante filtros.
\end{enumerate}

\subsection*{Variables requeridas por el dashboard}
\addcontentsline{toc}{subsection}{Variables requeridas por el dashboard}

Para cubrir los filtros y comparaciones solicitados fue necesario ampliar la ingeniería de variables del notebook. Se incorporaron \texttt{Income\_Segment} y \texttt{Spending\_Segment}, calculados mediante terciles de ingreso y cuartiles de gasto; \texttt{Dependent\_Composition} y \texttt{Marital\_Status\_Grouped}, para resumir la estructura familiar; \texttt{Dominant\_Channel} y \texttt{Channel\_Segment}, para distinguir el canal con más compras y el comportamiento mono o multicanal; y \texttt{Dominant\_Category}, que identifica la categoría con mayor gasto de cada cliente.

También se creó \texttt{Web\_Purchase\_Visit\_Ratio}, definido como compras web entre visitas web. Este indicador se utiliza únicamente de forma exploratoria y no se interpreta como una tasa de conversión real, porque las compras cubren dos años mientras que las visitas corresponden al último mes. Finalmente, las coordenadas \texttt{PCA\_PC1} y \texttt{PCA\_PC2} se incorporaron al archivo exportado para reutilizar en el dashboard la proyección ya calculada en el notebook.

Los segmentos estratégicos se representan mediante indicadores booleanos que pueden solaparse. Se consideran los perfiles \textit{Premium}, alto ingreso con bajo gasto, multicanal, web curioso, sensible a descuentos, leal, nuevo de alto valor e inactivo o frío. Un cliente puede pertenecer a más de una cohorte; por ello, estas etiquetas deben interpretarse como reglas descriptivas y no como clases mutuamente excluyentes.

\subsection*{Indicadores financieros y simulador}
\addcontentsline{toc}{subsection}{Indicadores financieros y simulador}

El resumen sin filtros reproduce los valores de la campaña original: 2{,}240 contactos, 334 respuestas, una tasa de aceptación de 14.91\%, un costo de 6{,}720 MU, un ingreso de 3{,}674 MU y un balance de $-3{,}046$ MU. Los cálculos utilizan las constantes \texttt{Z\_CostContact = 3} y \texttt{Z\_Revenue = 11}.

Cuando se aplican filtros, el dashboard recalcula dinámicamente el número de clientes seleccionados, las respuestas observadas, el costo, el ingreso, el balance y el ROI histórico. La última página compara este resultado con el contacto masivo y presenta los grupos priorizados por valor y aceptación. Esta operación responde a la pregunta ``¿qué habría ocurrido históricamente con este subconjunto?'', pero no constituye una predicción ni garantiza que la misma tasa de respuesta se repita en una campaña futura.

Para evitar fuga de información, \texttt{Response} no se ofrece como criterio de selección en el simulador. Sí pueden utilizarse variables disponibles antes de la campaña, como ingreso, gasto, edad, educación, antigüedad, recencia, canal, categoría, campañas anteriores aceptadas y quejas.

\subsection*{Diseño y validación}
\addcontentsline{toc}{subsection}{Diseño y validación}

La interfaz emplea una composición simple: navegación lateral, filtros contextuales, tarjetas KPI y gráficas interactivas construidas con Plotly. Se mantuvo una paleta reducida y fondos claros para priorizar legibilidad sobre ornamentación. Las siete páginas fueron ejecutadas y verificadas en el navegador; no se detectaron excepciones durante la carga, y se comprobó que los filtros actualizan las métricas y visualizaciones correspondientes.

\section*{Conclusiones}
\addcontentsline{toc}{section}{Conclusiones}

\textit{Pendiente de redacción después de validar las figuras y definir el alcance final del proyecto.}

\section*{Referencias}
\addcontentsline{toc}{section}{Referencias}

\textit{Pendiente de consolidar las fuentes en un formato bibliográfico uniforme.}

\end{document}
